\documentclass[11pt]{article}
\usepackage[margin=1in]{geometry}
\usepackage{amsmath,amsthm,amssymb}

\newtheorem*{finding}{Finding}
\newtheorem*{claim}{Claim}

\title{Analysis Notes: Error in NP-Hardness Proof and\\Polynomial-Time Algorithms for Fixed $\kappa$}
\date{April 2026}

\begin{document}
\maketitle

\noindent This document records the key findings from a careful review of the
companion documents on ORP complexity.

\section{Error in the NP-Hardness Proof}

\subsection{Location of the error}

The error is in \S3.3 of \texttt{np\_hardness\_proof\_revised.tex}, in the paragraph
titled ``Any permutation can be improved to a lifted permutation.''  The proof
asserts:

\begin{quote}
\emph{``By the realisability result, there exists a lifted permutation $\hat{\pi}^*$
--- constructed from some vertex ordering $\pi'$ --- that makes the same $M$
intra-block choices as $\hat{\pi}$.''}
\end{quote}

\noindent This is \textbf{false}.

\subsection{Why it is false}

A \emph{lifted} permutation corresponds to a linear ordering $\pi'$ of the $N$
tournament vertices.  The ``realisability result'' immediately above only establishes
that any binary choice vector $(c_1, \ldots, c_M) \in \{0,1\}^M$ yields a valid
permutation of the $2M$ events.  It says nothing about realisability as a
\emph{lifted} permutation, which additionally requires the $M$ binary choices to be
jointly consistent with a linear ordering.

For a cyclic tournament (e.g.\ $1 \to 2,\; 2 \to 3,\; 3 \to 1$), the choices
``agree with all three arcs simultaneously'' require
$\pi'(1) < \pi'(2) < \pi'(3) < \pi'(1)$, which is impossible.  No lifted permutation
achieves zero intra-block disagreements for a cyclic tournament.

\subsection{The actual minimum cost}

Despite the absence of any consistent lifted permutation for cyclic choices, the
\emph{non-lifted} permutation that processes blocks in slot order $1, 2, \ldots, M$
and within each block places events in their individually most-likely order achieves:
\begin{itemize}
\item zero intra-block disagreement cost (each block optimised independently), and
\item zero inter-block cost (blocks are listed in increasing slot order, so all
  inter-block pairs with $p = 1$ are respected).
\end{itemize}
Its total expected Kendall tau distance is therefore $M/4$ (only the fixed
$\min(p_r, 1-p_r) = 1/4$ terms) for \emph{every} tournament, regardless of cycle
structure.

Since the threshold is $K = k/2 + M/4 \ge M/4$ for any $k \ge 0$, the ORP instance
constructed by the reduction \emph{always} answers YES\@.  FAST does not always
answer YES.  The reduction is therefore incorrect.

\subsection{Why the construction is structurally broken}

The independent-block design makes the $2M$-event ORP problem decompose into $M$
independent binary decisions.  Each block can be solved greedily with no regard for
global consistency with a linear order.  The block structure that was intended to
isolate intra-block decisions is precisely what makes the problem trivially easy.

The pairwise decomposition of the Kendall tau distance (Lemma~2 of the main
document) provides no third-order terms that could penalise cyclic triples.  There is
therefore no mechanism within the ORP cost function to enforce transitivity across
blocks.

\subsection{The proof cannot be patched with this construction}

Patching would require the minimum ORP cost to equal $\mathrm{fas}(T)/2 + M/4$, but:
\begin{itemize}
\item The non-lifted optimum is $M/4$ for every tournament.
\item Making blocks interact would require cross-block posteriors encoding tournament
  structure, but there is no pairwise mechanism that penalises cyclic inconsistency.
\end{itemize}
Any correct reduction to ORP requires a fundamentally different construction, such as
$N$ coupled events (one per tournament vertex) whose joint posteriors encode the
desired tournament, or a reduction from a different NP-hard problem.

\section{What the Reduction \emph{Does} Establish}

The construction does correctly show two things:

\begin{enumerate}
\item \textbf{The $(\Rightarrow)$ direction is valid.}  If $\mathrm{fas}(T, \pi') \le
  k$ for some linear ordering $\pi'$, then the lifted permutation corresponding to
  $\pi'$ achieves cost $\le K$.  There is no error here.

\item \textbf{The per-block posteriors are computed correctly.}  The noise
  distributions in the construction give intra-block posteriors of $3/4$ or $1/4$
  as claimed, and inter-block pairs are indeed deterministic.
\end{enumerate}

Only the $(\Leftarrow)$ direction fails.

\section{Status of ORP Complexity}

\subsection{ORP with general distributions (part of input) is almost certainly NP-hard}

The ORP objective $\sum_{i<j} |2p_{ij}-1| \cdot \mathbf{1}[\hat\pi \text{ disagrees
on } (i,j)]$ is exactly the minimum weighted feedback arc set in a tournament
(MWFAS), where arc $(i,j)$ has weight $|2p_{ij}-1|$ and the direction is the
most-likely ordering.  MWFAS is NP-hard (it includes FAST as a special case with
uniform weights).

A correct NP-hardness proof would need to demonstrate that any FAST instance can be
encoded as an ORP instance (with $N$ events) whose minimum cost corresponds to
$\mathrm{fas}(T)$.  This requires showing that for any tournament $T$ on $N$ vertices,
there exist noise distributions $F_1, \ldots, F_N$ and timestamps $\hat{t}_1, \ldots,
\hat{t}_N$ such that the induced posteriors $p_{ij}$ point in the desired directions
with equal weights.  This is believed achievable (the product-form likelihood has
sufficient degrees of freedom) but requires a construction that has not yet been
worked out rigorously.

\subsection{ORP with fixed $\kappa$ and MLR distributions is polynomial}

The companion document \texttt{general\_kappa\_algorithm.tex} establishes:

\begin{claim}
For any fixed catalogue size $\kappa$ and noise distributions satisfying the
monotone likelihood ratio (MLR) property (equivalently, log-concavity of the density
--- satisfied by Gaussian, Laplacian, logistic, uniform on an interval, geometric,
Poisson, and all other standard unimodal distributions), ORP is solvable in
$O(N^\kappa)$ time.
\end{claim}

The algorithm is a $\kappa$-dimensional dynamic programming interleaving algorithm:
\begin{enumerate}
\item Within each same-type group, the optimal ordering is the timestamp ordering
  (Lemma~1 of \texttt{general\_kappa\_algorithm.tex}: MLR implies $p_{ij} \ge 1/2
  \iff \hat{t}_i \le \hat{t}_j$ for same-type pairs).
\item Sort each group by timestamp; this reduces the problem to finding the optimal
  interleaving of $\kappa$ sorted sequences.
\item A $\kappa$-dimensional DP with $O(N^\kappa)$ states and $O(\kappa)$-time
  transitions (using precomputed prefix sums) finds the optimal interleaving.
\end{enumerate}

\noindent The complexity table is:

\medskip
\begin{center}
\renewcommand{\arraystretch}{1.3}
\begin{tabular}{@{}cll@{}}
\hline
\textbf{Catalogue size} & \textbf{Complexity (MLR distributions)} & \textbf{Algorithm} \\
\hline
$\kappa = 1$ & $O(N \log N)$ & Sort by timestamp \\
$\kappa = 2$ & $O(N^2)$ & 2-D interleaving DP \\
$\kappa = 3$ & $O(N^3)$ & 3-D interleaving DP \\
Fixed $\kappa$ & $O(N^\kappa)$ & $\kappa$-D interleaving DP \\
\hline
\end{tabular}
\end{center}

\medskip
\noindent Note that the $\kappa = 3$ case is \textbf{polynomial} --- the claim in
\texttt{two\_type\_algorithm.tex} that $\kappa \ge 3$ is NP-hard (for MLR
distributions) is not established by the existing proof.

\subsection{Non-MLR distributions}

Lemma~1 of \texttt{general\_kappa\_algorithm.tex} can fail for distributions that
violate MLR\@.  A concrete example: a bimodal distribution $G$ with mass at
$\{-3, 0, 3\}$ concentrated on $\pm 3$ can make the most-likely ordering for a
same-type pair \emph{disagree} with the timestamp ordering, since the data is more
consistent with the two events each having large noise of opposite sign than with
both having near-zero noise.  For such distributions, the within-group ordering
problem is non-trivial, and NP-hardness may be achievable via a correct reduction.

\section{Recommended Follow-up}

\begin{enumerate}
\item \textbf{Retract or annotate} the NP-hardness claim in
  \texttt{np\_hardness\_proof\_revised.tex} and \texttt{proof\_summary.tex} until a
  correct proof is found.
\item \textbf{Determine} whether any FAST instance can be encoded as an $N$-event ORP
  instance with posteriors in the desired directions.  If yes, write a corrected
  reduction.
\item \textbf{Determine} whether ORP with non-MLR distributions and fixed $\kappa$ is
  NP-hard.  The bimodal distribution example suggests that non-MLR distributions
  could enable a valid reduction even for $\kappa = 3$.
\end{enumerate}

\end{document}
